\section{Edit Distance and Sequence Alignment}\label{dp: sec: edit distance}\doHeader{Edit Distance}

\begin{mylist}
\item[\bf input:]
  A pair $I=(A[1..m], B[1..n])$ of strings.

\item[\bf output:]
  The \emph{edit distance} between $A$ and $B$.
  This is the minimum cost of any \emph{alignment}
  between $A$ and $B$.
  An alignment is a matching between characters of $A$
  and characters of $B$ such that the matching doesn't \emph{cross}:
  that is, there are no two matched pairs $(A[i], B[j])$ and $(A[i'], B[j'])$
  where $i < i'$ but $j > j'$.
  The cost of an alignment is the total number of unmatched characters (in both sides)
  plus the number of matched pairs $(A[i], B[j])$ where $A[i] \ne B[j]$.
\end{mylist}
Equivalently, the edit distance is the minimum number
of edit operations that transforms $A$ into $B$,
where each edit operation either deletes a character,
inserts a character, or replaces any one character by some other character.

\paragraph{intuition.}
Fix any input instance $(A[1..m], B[1..n])$.
The collection of objects we are working with is the set of alignments between $A$ and $B$.
To derive the algorithm, we classify these alignments into three types:
\begin{mylist}
\item[(A)] Those that don't match $A[m]$ at all.
\item[(B)] Those that  don't match $B[n]$ at all. 
\item[(C)] Those that match $A[m]$ to $B[n]$.
\end{mylist}
Every alignment is of one of these types, because an alignment cannot both match $A[m]$
to some character in $B[1..n-1]$ and match $B[n]$ to some character in $A[1..m-1]$.
(This is because in such a matching those matched pairs would cross.)
It is possible for an alignment to be both type (A) and type (B),
but this is not a problem for the algorithm.
Next we consider the three types of alignments,
and for each we observe that the problem of finding a minimum-cost alignment of that type
reduces to a smaller subproblem:
\begin{mylist}
% [review 2026-09-27] fix: m and n were swapped (input is A[1..m], B[1..n]); type (A) now uses A[1..m-1] and B[1..n]
\item[(A)] An alignment of type (A) is just an alignment of $A[1..m-1]$ and $B[1..n]$,
  so the minimum cost of any alignment of type (A)
  is just the minimum cost of any alignment between $A[1..m-1]$ and $B[1..n]$.

% [review 2026-09-27] fix: m and n were swapped (input is A[1..m], B[1..n]); type (B) now uses A[1..m] and B[1..n-1]
\item[(B)] Likewise, the minimum cost of any alignment of type (B)
  is just the minimum cost of any alignment between $A[1..m]$ and $B[1..n-1]$.

% [review 2026-09-27] fix: m and n were swapped (input is A[1..m], B[1..n]); type (C) now matches A[m] to B[n] and uses A[1..m-1], B[1..n-1]
\item[(C)] Any alignment of type (C) matches $A[m]$ to $B[n]$, and then the rest of the alignment
  is just an alignment between $A[1..m-1]$ and $B[1..n-1]$.
  So, the minimum cost of any alignment of type (C)
  is just the minimum cost of any alignment between  $A[1..m-1]$ and $B[1..n-1]$,
  plus the cost of matching $A[m]$ to $B[n]$
  (1 if $A[m] \ne B[n]$, otherwise 0).
\end{mylist}

When we apply this reasoning recursively,
each subproblem that arises is to compute
the minimum cost of any alignment between some prefix of $A$
and some prefix of $B$.
That is, between $A[1..i]$ and $B[1..j]$ for some $i$ and $j$.
Here is the algorithm:

\paragraph{subproblems and recurrence relation.}
For each pair $(i, j)$ with $i\in\{0,1,\ldots, m\}$ and $j\in\{0,1,\ldots,n\}$,
define $M(i, j)$ to be the minimum cost of any alignment between $A[1..i]$ and $B[1..j]$.
% [review 2026-09-27] fix: final answer was M(n,m) but i ranges over 0..m and j over 0..n; changed to M(m,n)
The final answer is $M(m, n)$.  The recurrence relation is
\[
  M(i, j) =
  \begin{cases}
    \max(i, j) & \text{ if } i=0 \text{ or } j=0, \\
    \min (M(i-1, j) + 1, M(i, j-1) + 1, M(i-1, j-1)) & \text{ if } i,j > 0 \text{ and } A[i] = B[j],  \\
    \min (M(i-1, j) + 1, M(i, j-1) + 1, M(i-1, j-1) + 1) & \text{ if } i,j > 0 \text{ and } A[i] \ne B[j]. 
  \end{cases}
\]

\paragraph{correctness.}
Correctness follows from the ideas sketched above under ``intuition.''

\paragraph{time.}
There are $(n+1)\cdot(m+1) = O(nm)$ subproblems,
and for each, the right-hand side of the recurrence takes
constant time to evaluate
(assuming the smaller subproblems are already solved).
So the algorithm (iterative, or recursive and memoized)
take $O(nm)$.

\begin{theorem}
  There is an $O(nm)$-time algorithm for \prob{Edit Distance}.
\end{theorem}

\pythonCode{dynamic_programming/edit_distance.py}

\subsection*{Extending the algorithm to Sequence Alignment}\doHeader{Sequence Alignment}

\KT Section 6.6 considers a generalization of \prob{Edit Distance} called \prob{Sequence Alignment}:

\begin{mylist}
\item[\bf input:]
  A pair $I=(A[1..m], B[1..n])$ of strings, \emph{gap penalty} $\delta>0$,
  and \emph{mismatch cost function} $\alpha$ defining,
  for each character $a$ in $A$ and $b$ in $B$,
  a \emph{mismatch cost} $\alpha(a, b)\ge 0$.

\item[\bf output:]
  The minimum cost of any alignment of $A$ and $B$,
  where an alignment is a non-crossing matching (as defined for \prob{Edit Distance}),
  but the cost of the alignment is $\delta$ times the total number of unmatched characters (in both sides),
  plus the sum over the matched pairs $(A[i], B[j])$ of $\alpha(A[i], B[j])$.
\end{mylist}
\prob{Edit Distance} is the special case when $\delta=1$
and $\alpha(a, b) = 1$ for $a\ne b$, and $\alpha(a, b) = 0$ for $a=b$.
Following the same reasoning as for the algorithm for \prob{Edit Distance}
gives the following algorithm.
Fix any input instance $(A[1..m], B[1..n])$.

\paragraph{subproblems and recurrence relation.}
For each pair $(i, j)$ with $i\in\{0,1,\ldots, m\}$ and $j\in\{0,1,\ldots,n\}$,
define $M(i, j)$ to be the minimum cost of any alignment between $A[1..i]$ and $B[1..j]$.
% [review 2026-09-27] fix: final answer was M(n,m) but i ranges over 0..m and j over 0..n; changed to M(m,n)
The final answer is $M(m, n)$.  The recurrence relation is
\[
  M(i, j) = \begin{cases}
    \delta \max(i, j) & \text{ if } i=0 \text{ or } j=0 \\
    \min \big(M(i-1, j) + \delta, M(i, j-1) + \delta, M(i-1, j-1) + \alpha(A[i], B[j])\big)  & \text{ if } i, j > 0.
  \end{cases}
\]

\paragraph{correctness.}  Correctness follows by reasoning similar to the reasoning
for the correctness of the recurrence relation for \prob{Edit Distance}.

\paragraph{time.}
There are $(n+1)\cdot(m+1) = O(nm)$ subproblems,
and for each, the right-hand side of the recurrence takes
constant time to evaluate
(assuming the smaller subproblems are already solved).
So the algorithm (iterative, or recursive and memoized) takes time $\Theta(nm)$.

\begin{theorem}
  There is an $O(nm)$-time algorithm for \prob{Sequence Alignment}.
\end{theorem}

% \paragraph{External sources on \prob{Edit Distance}}
  
% \providecommand{\URL}[1]{\parbox[t]{0.8\textwidth}{\url{#1}}}

% \begin{itemize}
% \item \JE Section 3.7. KT 6.6. CLRS Problem 15--5 (Chapter 15). DPV~Chapter 6.3. MIT video:

% \URL{https://ocw.mit.edu/courses/electrical-engineering-and-computer-science/6-006-introduction-to-algorithms-fall-2011/lecture-videos/lecture-21-dp-iii-parenthesization-edit-distance-knapsack}

% \end{itemize}

%%% Local Variables:
%%% mode: latex
%%% TeX-master: "dynamic_programming"
%%% End:
